\documentclass[12pt]{article}
%^ Start of document
\usepackage{times}  %Makes text TimesNewRoman
\usepackage{setspace} %Allows you to set single or double space 
\usepackage{biblatex} %Imports biblatex package
\usepackage{graphicx} % Required for inserting images
\usepackage{authblk}
\usepackage[colorlinks=true, urlcolor=blue, linkcolor=blue]{hyperref}
\usepackage{subfig}
\usepackage{float}
\usepackage{tabularx}
\usepackage{multirow}
\usepackage[paperheight=11in,
   paperwidth=8.5in,
   top=20mm,
   bottom=20mm,
   left=40mm,
   right=40mm]{geometry}
\usepackage{moresize}
\usepackage{tikz}
\usepackage{xcolor}
\usepackage{physics}
\usepackage{amssymb}
\usepackage{mathrsfs}
\usepackage{amsmath}
\usepackage{ragged2e}
\usepackage{comment}
\usepackage{xcolor}
\addbibresource{mathmeth.bib}
\begin{document}
\singlespacing  %Sets single spacing for whole document
\title{Math Methods HW 3}

\author[1]{Zachary Tullier}
\affil[1]{Student\\Physics Department\\ University of Houston - Clear Lake \\Houston, Texas\\U.S}
\date{}
\maketitle




\section*{Textbook Question 10.1.5}
    Construct the Green's function for 
    \[
        x^2 \frac{d^2y}{dx^2} + x \frac{dy}{dx} + (k^2 x^2 - 1)y = 0
    \]
    subject to the boundary conditions $y(0) = 0$, $y(1) = 0$.

    \subsection*{Solution:}
        The given equation is a second-order linear ordinary differential equation (ODE) of the form:
        \[
        \mathcal{L}y = 0,
        \]
        where the operator \(\mathcal{L}\) is defined as:
        \[
        \mathcal{L}y = x^2 \frac{d^2 y}{dx^2} + x \frac{dy}{dx} + \left(k^2 x^2 - 1\right)y.
        \]
        The Green's function, \(G(x, \xi)\), satisfies:
        \[
        \mathcal{L}G(x, \xi) = \delta(x - \xi),
        \]
        with the same boundary conditions:
        \[
        G(0, \xi) = 0, \quad G(1, \xi) = 0.
        \]
        We solve the homogeneous equation:
        \[
        x^2 \frac{d^2 y}{dx^2} + x \frac{dy}{dx} + \left(k^2 x^2 - 1\right)y = 0.
        \]
        This is a Cauchy-Euler equation, and its solutions are generally expressed in terms of Bessel functions. Let:
        \[
        y(x) = x^m.
        \]
        Substitute \(y(x) = x^m\) into the equation:
        \[
        x^2 \cdot m(m-1)x^{m-2} + x \cdot mx^{m-1} + \left(k^2 x^2 - 1\right)x^m = 0.
        \]
        Simplify:
        \[
        m(m-1) + m + (k^2 x^2 - 1) = 0.
        \]
        This leads to solutions involving Bessel functions. The general solution is:
        \[
        y(x) = A J_1(kx) + B Y_1(kx),
        \]
        where \(J_1(kx)\) and \(Y_1(kx)\) are the Bessel functions of the first and second kind, respectively.
        The Green's function, \(G(x, \xi)\), is constructed using the two linearly independent solutions \(J_1(kx)\) and \(Y_1(kx)\). We write:
        \[
        G(x, \xi) =
        \begin{cases} 
        A_1 J_1(kx) + B_1 Y_1(kx), & x < \xi, \\
        A_2 J_1(kx) + B_2 Y_1(kx), & x > \xi.
        \end{cases}
        \]
        The constants \(A_1, B_1, A_2, B_2\) are determined by the following conditions:
        \begin{enumerate}
            \item Continuity of \(G(x, \xi)\) at \(x = \xi\):**
                \[
                    G(\xi^-, \xi) = G(\xi^+, \xi).
                \]
            \item Discontinuity in the first derivative at \(x = \xi\):**
                \[
                    \frac{dG}{dx}\Big|_{x=\xi^+} - \frac{dG}{dx}\Big|_{x=\xi^-} = \frac{1}{\xi}.
                \]
            \item Boundary conditions:
                \[
                    G(0, \xi) = 0, \quad G(1, \xi) = 0.
                \]
        \end{enumerate}
        Apply boundary condition, \(G(0, \xi) = 0\), and \(Y_1(kx)\) diverges at \(x = 0\), we must set \(B_1 = 0\). Thus:
        \[
        G(x, \xi) =
        \begin{cases} 
        A_1 J_1(kx), & x < \xi, \\
        A_2 J_1(kx) + B_2 Y_1(kx), & x > \xi.
        \end{cases}
        \]
        Apply boundary condition, \(G(1, \xi) = 0\), we have:
        \[
        A_2 J_1(k) + B_2 Y_1(k) = 0.
        \]
        Solve for \(B_2\) in terms of \(A_2\):
        \[
        B_2 = -A_2 \frac{J_1(k)}{Y_1(k)}.
        \]
        Continuity at \(x = \xi\) implies
        \[
        A_1 J_1(k\xi) = A_2 J_1(k\xi) + B_2 Y_1(k\xi).
        \]
        At \(x = \xi\), the derivative discontinuity condition is:
        \[
        \frac{dG}{dx}\Big|_{x=\xi^+} - \frac{dG}{dx}\Big|_{x=\xi^-} = \frac{1}{\xi}.
        \]
        For \(x < \xi\):
        \[
        \frac{dG}{dx} = A_1 k J_1'(kx).
        \]
        For \(x > \xi\):
        \[
        \frac{dG}{dx} = A_2 k J_1'(kx) + B_2 k Y_1'(kx).
        \]
        At \(x = \xi\), this becomes:
        \[
        A_2 k J_1'(k\xi) + B_2 k Y_1'(k\xi) - A_1 k J_1'(k\xi) = \frac{1}{\xi}.
        \]
        Solve Equations to determine the constants \(A_1, A_2, B_2\). Substitute them into the Green's function expression:
        \[
        G(x, \xi) =
        \begin{cases} 
        A_1 J_1(kx), & x < \xi, \\
        A_2 J_1(kx) + B_2 Y_1(kx), & x > \xi.
        \end{cases}
        \]
        The detailed expression for \(G(x, \xi)\) depends on solving these constants explicitly. The final Green's function involves the Bessel functions \(J_1(kx)\) and \(Y_1(kx)\), ensuring the boundary conditions and continuity properties are satisfied.

\section*{Textbook Question 10.2.2}
    Show that if 
    \[
        \mathcal{L} \psi(r) \equiv \nabla \cdot \left( p(r) \nabla \psi(r) \right) + q(r) \psi(r)
    \]
    then $\mathcal{L}$ is hermitian for $p(r)$ and $q(r)$ real, assuming Dirichlet boundary conditions on the boundary of a region and that the scalar product is an integral over that region with unit weight.
   
   \subsection*{Solution:}
        An operator \(\mathcal{L}\) is Hermitian if, for any two functions \(\psi(\mathbf{r})\) and \(\phi(\mathbf{r})\), the scalar product satisfies:
        \[
        \langle \phi, \mathcal{L}\psi \rangle = \langle \mathcal{L}\phi, \psi \rangle,
        \]
        where the scalar product is defined as:
        \[
        \langle \phi, \psi \rangle = \int_V \phi^*(\mathbf{r}) \psi(\mathbf{r}) \, d^3r,
        \]
        and \(V\) is the volume of the region of interest.
        Since Dirichlet boundary conditions are assumed, \(\psi(\mathbf{r}) = 0\) and \(\phi(\mathbf{r}) = 0\) on the boundary of the region.
        Expand \(\langle \phi, \mathcal{L}\psi \rangle\)
        Using the definition of \(\mathcal{L}\), compute \(\langle \phi, \mathcal{L}\psi \rangle\):
        \[
        \langle \phi, \mathcal{L}\psi \rangle = \int_V \phi^*(\mathbf{r}) \left[ \nabla \cdot \left( p(\mathbf{r}) \nabla \psi(\mathbf{r}) \right) + q(\mathbf{r}) \psi(\mathbf{r}) \right] \, d^3r.
        \]
        Split this integral into two terms:
        \[
        \langle \phi, \mathcal{L}\psi \rangle = \int_V \phi^*(\mathbf{r}) \nabla \cdot \left( p(\mathbf{r}) \nabla \psi(\mathbf{r}) \right) \, d^3r + \int_V \phi^*(\mathbf{r}) q(\mathbf{r}) \psi(\mathbf{r}) \, d^3r.
        \]
        First Term - Divergence Integral
        Apply the divergence theorem to the first term:
        \[
        \int_V \phi^*(\mathbf{r}) \nabla \cdot \left( p(\mathbf{r}) \nabla \psi(\mathbf{r}) \right) \, d^3r = \int_S \phi^*(\mathbf{r}) p(\mathbf{r}) \nabla \psi(\mathbf{r}) \cdot \mathbf{n} \, dS - \int_V \nabla \phi^*(\mathbf{r}) \cdot \left( p(\mathbf{r}) \nabla \psi(\mathbf{r}) \right) \, d^3r,
        \]
        where \(S\) is the boundary of the region, and \(\mathbf{n}\) is the outward unit normal vector.
        Under Dirichlet boundary conditions (\(\psi = 0\) and \(\phi = 0\) on the boundary), the surface integral vanishes:
        \[
        \int_S \phi^*(\mathbf{r}) p(\mathbf{r}) \nabla \psi(\mathbf{r}) \cdot \mathbf{n} \, dS = 0.
        \]
        Thus:
        \[
        \int_V \phi^*(\mathbf{r}) \nabla \cdot \left( p(\mathbf{r}) \nabla \psi(\mathbf{r}) \right) \, d^3r = -\int_V \nabla \phi^*(\mathbf{r}) \cdot \left( p(\mathbf{r}) \nabla \psi(\mathbf{r}) \right) \, d^3r.
        \]
        Expand the dot product:
        \[
        -\int_V \nabla \phi^*(\mathbf{r}) \cdot \left( p(\mathbf{r}) \nabla \psi(\mathbf{r}) \right) \, d^3r = -\int_V p(\mathbf{r}) \nabla \phi^*(\mathbf{r}) \cdot \nabla \psi(\mathbf{r}) \, d^3r.
        \]
        Second Term - Potential Term
        The second term is:
        \[
        \int_V \phi^*(\mathbf{r}) q(\mathbf{r}) \psi(\mathbf{r}) \, d^3r.
        \]
        Combine Terms for \(\langle \phi, \mathcal{L}\psi \rangle\)
        Combine the results for the two terms:
        \[
        \langle \phi, \mathcal{L}\psi \rangle = -\int_V p(\mathbf{r}) \nabla \phi^*(\mathbf{r}) \cdot \nabla \psi(\mathbf{r}) \, d^3r + \int_V \phi^*(\mathbf{r}) q(\mathbf{r}) \psi(\mathbf{r}) \, d^3r.
        \]
        Expand \(\langle \mathcal{L}\phi, \psi \rangle\)
        Similarly, compute \(\langle \mathcal{L}\phi, \psi \rangle\):
        \[
        \langle \mathcal{L}\phi, \psi \rangle = \int_V \psi^*(\mathbf{r}) \left[ \nabla \cdot \left( p(\mathbf{r}) \nabla \phi(\mathbf{r}) \right) + q(\mathbf{r}) \phi(\mathbf{r}) \right] \, d^3r.
        \]
        Following the same steps as above, using the divergence theorem and Dirichlet boundary conditions:
        \[
        \langle \mathcal{L}\phi, \psi \rangle = -\int_V p(\mathbf{r}) \nabla \psi^*(\mathbf{r}) \cdot \nabla \phi(\mathbf{r}) \, d^3r + \int_V \psi^*(\mathbf{r}) q(\mathbf{r}) \phi(\mathbf{r}) \, d^3r.
        \]
        Symmetry and Hermiticity
        Compare \(\langle \phi, \mathcal{L}\psi \rangle\) and \(\langle \mathcal{L}\phi, \psi \rangle\):
        \[
        \langle \phi, \mathcal{L}\psi \rangle = -\int_V p(\mathbf{r}) \nabla \phi^*(\mathbf{r}) \cdot \nabla \psi(\mathbf{r}) \, d^3r + \int_V \phi^*(\mathbf{r}) q(\mathbf{r}) \psi(\mathbf{r}) \, d^3r,
        \]
        \[
        \langle \mathcal{L}\phi, \psi \rangle = -\int_V p(\mathbf{r}) \nabla \psi^*(\mathbf{r}) \cdot \nabla \phi(\mathbf{r}) \, d^3r + \int_V \psi^*(\mathbf{r}) q(\mathbf{r}) \phi(\mathbf{r}) \, d^3r.
        \]
        Since \(p(\mathbf{r})\) and \(q(\mathbf{r})\) are real, and the scalar product involves complex conjugates, these two expressions are equal:
        \[
        \langle \phi, \mathcal{L}\psi \rangle = \langle \mathcal{L}\phi, \psi \rangle.
        \]
        The operator \(\mathcal{L}\) is Hermitian under Dirichlet boundary conditions with real \(p(\mathbf{r})\) and \(q(\mathbf{r})\).

\end{document}
